\cotefile{63}
\scanpage{78}
\mspage{\handwritten{- 1 -}}

\noindent\uline{A\,V\,A\,N\,T\,-\,P\,R\,O\,P\,O\,S}

\medskip
\noindent\begin{minipage}{\linewidth}\small
I.\qquad
% One continuous oblique pen stroke, from above the opening down into the footnote.
\tikz[remember picture,overlay]{\coordinate (cotesixtythreestart) at (-0.7em,1.8ex);}
Le but principal du présent Séminaire est de développer le formalisme de la ``cohomologie de Weil'' des Schémas. A partir essentiellement des résultats qui sont démontrés ici, des arguments bien connus, d'ailleurs dûs à Weil lui-même, \unsure{\'{p}|p}ermettent de déduire une partie des conjectures de Weil sur les fonctions $\mathrm{L}$ des variétés projectives non singulières sur un corps fini\textsuperscript{\handwritten{(1)}}. Cette déduction sera d'ailleurs présentée dans SGA \handwritten{1964/65}, dans un séminaire qui fera suite à celui-ci.

Dans le présent séminaire, nous nous bornons à l'étude de la cohomologie des schémas relativement à la \uline{topologie étale} . Cette topologie par sa description est fort proche des topologies des variétés topologiques habituelles\handwritten{,} et on verra que la plupart des résultats classiques concernant la cohomologie des espaces topologiques\unsure{.}ordinaires (suites spectrales variées, théorèmes de finitude, K\handwritten{\"}{u}nneth, dualité, théorème de Lefschetz) peuvent se formuler et se démontrer dans le nouveau contexte, à condition de se borner le cas échéant aux faisceaux de torsion \uline{premiers aux caractéristiques résiduelles} des schémas envisagés\unsure{.|,} On obtiendra une théorie cohomologique ``à coefficient de caractéristique $0$'' (comme demandée par Weil) par un passage à la limite projective e\makebox[0pt][l]{\deleted{\unsure{ss|xx|ee}}}ssentiellement trivial (cf SGA \handwritten{1965}), permettant de définir une cohomologie à coefficient d\deleted{\unsure{e}}\handwritten{\unsure{a|e}ns} l'anneau $\underline{\mathrm{Z}_{\mathord{\rlap{\handwritten{$\scriptstyle/$}}\unsure{\mathrm{l}|1}}}}$ des entiers $\mathord{\rlap{\handwritten{$/$}}\mathrm{l}}$-adiques à partir des coefficients $\underline{\mathrm{Z}}\handwritten{/}\mathord{\rlap{\handwritten{$/$}}\mathrm{l}}^{\handwritten{\unsure{\nu|v}}}\underline{\mathrm{Z}}$ \unsure{.}, $\handwritten{\nu\longrightarrow}+\handwritten{\infty}$ . Lorsque $\mathord{\rlap{\handwritten{$/$}}\mathrm{l}}$ est premier aux caractéristiques résiduelles, cette cohomologie possède toutes les bonnes propriétés habituelles dans la cohomologie à coefficient $\underline{\mathrm{Z}}$ classique (et se prête donc à la formulation des conjectures de Weil). D'ailleurs, lorsque les schémas envisagés sont de type fini sur le corps des complexes $\underline{\mathrm{C}}$, nous verrons que la cohomologie (à coefficient\handwritten{s} de torsion) est essentiellement identique à la cohomologie

\smallskip
\noindent\rule{0.32\linewidth}{0.4pt}\par
{\footnotesize\noindent\textsuperscript{\handwritten{(1)}}\quad Au moment d'écrire ces lignes, il n'est pas prouvé que les valeurs propres de l'homomorphisme de Frobenius opérant sur les $\mathrm{H}^{\handwritten{i}}\handwritten{(}\mathrm{X},\underline{\mathrm{Z}}_{\handwritten{\ell}}\handwritten{)}$ sont des entiers algébriques, ni à fortiori que \begingroup\let\sout\mbox\deleted{leurs\handwritten{\tikz[remember picture,overlay]{\draw[line width=0.4pt] (cotesixtythreestart) -- (-0.45em,0.15ex);}}}\endgroup{} valeurs absolues sont égales à $\mathrm{q}^{\handwritten{i/2}}$ .\par}
\end{minipage}
% archive sheet number: 77
\clearpage
\scanpage{80}
% Short oblique impression at the upper left of the sheet.
\reversemarginpar
\margin{\figure{\tikz[x=1em,y=1em]{\draw[line width=.23pt] (0,0) -- (1.5,1.25); \fill (0,0) -- (.15,.02) -- (.25,.19) -- (.14,.19) -- cycle; \fill (1.37,1.11) -- (1.47,1.27) -- (1.53,1.22) -- (1.43,1.13) -- cycle;}}}
\normalmarginpar
\mspage{X-1}

\begin{center}
Liste des courbes elliptiques à multiplication complexe\\
dont l'invariant modulaire $\mathrm{j}$ est rationnel

\medskip
-{}-{}-
\end{center}

Ces courbes s'obtiennent à partir des ordres de corps quadratiques imaginaires de nombre de classes $1$. On doit donc :

a) faire la liste des corps quadratiques imaginaires $\underline{\underline{\mathrm{Q}}}(\handwritten{\sqrt}{-\mathrm{d}})$ dont l'ordre maximal (i.e. l'anneau des entiers) est principal. Le résultat est :
\reversemarginpar
\margin{\handwritten{\uline{NB}\quad les $\mathrm{d}$\\sont premiers !}}
\normalmarginpar
\[
\begin{aligned}
&\mathrm{d} = \handwritten{\overbrace}{1,\,2,\,3,\,7,\,11}^{\inserted{\handwritten{\text{``\unsure{euclidiens|euclidien}''}}}},\,19,\,43,\,67,\,163 \quad \text{plus peut-être}\\
&\text{une autre valeur}\quad \mathrm{d}_{?}\quad .
\end{aligned}
\]

b) Pour chacune des valeurs de $\mathrm{d}$ ci-dessus, déterminer quels sont les conducteurs $\underline{\mathrm{f}}$ qui conduisent à un ordre ayant un nombre de classes égal à $1$.

On trouve $\underline{\mathrm{f}}=1,2$ pour $\mathrm{d}=1$ , $\underline{\mathrm{f}}=1,2,3$ pour $\mathrm{d}=3$, $\underline{\mathrm{f}}=1,2$ pour $\mathrm{d}=7$, et $\underline{\mathrm{f}}=1$ pour $\mathrm{d}=2,\,11,\,19,\,43,\,67,\,163$ et $\mathrm{d}_{?}$ (s'il existe).

c) Pour chaque couple $(\mathrm{d},\underline{\mathrm{f}})$, prendre une base $(1,\mathrm{t})$ de l'ordre correspondant, et calculer $\mathrm{j}(\mathrm{t})$; c'est l'invariant modulaire cherché; il est rationnel.

Les valeurs de $\mathrm{j}(\mathrm{t})$ se trouvent toutes dans H.WEBER, \uline{Lehrbuch der Algebra}, t.III; de façon plus précise, Weber donne les valeurs de deux fonctions $\mathrm{f}(\mathrm{t})$ et $\mathrm{f}_{1}(\mathrm{t})$ liées à $\mathrm{j}$ par lesrelations
\[
\mathrm{j}=(\mathrm{f}^{24}-16)^{3}/\mathrm{f}^{24}
       =(\mathrm{f}_{1}^{24}+16)^{3}/\mathrm{f}_{1}^{24}\quad .
\]

On obtient ainsi les résultats suivants :

\begingroup
\makeatletter\tagsleft@true\makeatother
\[
(\mathrm{d}=1,\ \underline{\mathrm{f}}=1)\quad \rightarrow\quad
\mathrm{j}=\mathrm{j}(\mathrm{i})=2^{6}3^{3}
\tag*{1.1.}
\]
(courbe elliptique ayant un groupe d'automorphismes d'ordre $4$)

\[
(\mathrm{d}=1,\ \underline{\mathrm{f}}=2)\quad \rightarrow\quad
\mathrm{j}=\mathrm{j}(2\mathrm{i})=(2.3.11)^{3}
\tag*{1.2}
\]
\endgroup
% archive sheet number: 79
\clearpage
\scanpage{81}
\reversemarginpar
\margin{\figure{\tikz[x=1em,y=1em]{\draw[black!18,line width=.4pt] (0,0) -- (1.5,1.25); \fill (0,0) -- (.14,.02) -- (.22,.18) -- (.13,.20) -- cycle; \fill (1.34,1.08) -- (1.48,1.28) -- (1.54,1.24) -- (1.44,1.09) -- cycle;}}}
\normalmarginpar
\mspage{X-2}

\diagram{\[
\text{\small
\setlength{\arraycolsep}{.25em}
\resizebox{\linewidth}{!}{$\begin{array}{@{}r l c l l@{}}
2.1 & (\mathrm{d}=2,\ \underline{\mathrm{f}}=1) & \rightarrow &
\mathrm{j}=\mathrm{j}(\handwritten{\sqrt}{-2})=(2^{2}.5)^{3} & \text{(Weber, p.721)}\\[.6ex]
3.1 & (\mathrm{d}=3,\ \underline{\mathrm{f}}=1) & \rightarrow &
\mathrm{j}=\mathrm{j}(\dfrac{-1+\handwritten{\sqrt}{-3}}{2})=0 & \\[.3ex]
& \multicolumn{4}{l}{\quad\text{(courbe elliptique ayant un groupe d'autom.d'ordre }6\text{)}}\\[.6ex]
3.2 & (\mathrm{d}=3,\ \underline{\mathrm{f}}=2) & \rightarrow &
\mathrm{j}=\mathrm{j}(\handwritten{\sqrt}{-3})=2^{4}3^{3}5^{3} & \text{(Weber, p.721)}\\[.6ex]
3.3 & (\mathrm{d}=\rlap{$\deleted{\unsure{3|6}}$}\handwritten{\unsure{3|6}},\ \underline{\mathrm{f}}=3\handwritten{)} & \rightarrow &
\mathrm{j}=\mathrm{j}(\dfrac{-1+3\handwritten{\sqrt}{-3}}{2})=-3.2^{15}5^{3} & \text{(Weber,462)}\\[.6ex]
7.1 & (\mathrm{d}=7,\ \underline{\mathrm{f}}=1) & \rightarrow &
\mathrm{j}=\mathrm{j}(\dfrac{-1+\handwritten{\sqrt}{-7}}{2})=-3^{3}5^{3} & \text{(Weber,460)}\\[.6ex]
7\unsure{.|,}2 & (\mathrm{d}=7,\ \underline{\mathrm{f}}=2) & \rightarrow &
\mathrm{j}=\mathrm{j}(\handwritten{\sqrt}{-7})=(3.5.17)^{3} & \text{(Weber, 475)}\\[.6ex]
11.1 & (\mathrm{d}=11,\ \underline{\mathrm{f}}=1) & \rightarrow &
\mathrm{j}=\mathrm{j}\bigl(\dfrac{-1+\handwritten{\sqrt}{-1\rlap{$1$}}\phantom{1}}{2}\bigr)=-2^{15} & \text{(Weber, 462)}\\[.6ex]
19.1 & & \rightarrow & \mathrm{j}=-(2^{5}.3)^{3} & \multicolumn{1}{c}{\text{''}}\\[.6ex]
43.1 & & \rightarrow\,\raisebox{.5ex}{\scriptsize\unsure{,|.}} & \mathrm{j}=-(2^{6}.3.5)^{3} & \multicolumn{1}{c}{\text{''}}\\[.6ex]
67.1 & & \rightarrow & \mathrm{j}=-(2^{5}.3.5.11)^{3} & \multicolumn{1}{c}{\text{''}}\\[.6ex]
163.1 & & \rightarrow & \mathrm{j}=-(2^{6}.3.5.23.29)^{3} & \multicolumn{1}{c}{\text{''}}\\[.6ex]
\mathrm{d}_{?}.1 & & \rightarrow & \mathrm{j}=\ ? &
\end{array}$}}
\]}

\begin{center}
***
\end{center}
% archive sheet number: 80
\clearpage
\scanpage{83}
\reversemarginpar
\margin{\figure{\tikz[x=1em,y=1em]{\draw[line width=.22pt] (0,0) -- (1.48,-.04); \draw[black!48,line width=.2pt] (1.48,-.04) -- (2.02,-.05); \fill (0,0) ellipse (.13 and .055); \fill (1.48,-.03) ellipse (.045 and .033); \draw (2,-.01) -- (2.03,-.09); \draw[black!50,line width=.26pt] (.40,.64) .. controls (.51,.61) and (.50,.58) .. (.43,.53); \draw[black!22,line width=.18pt] (.44,.51) -- (.40,.04);}}}
\normalmarginpar
\mspage{B-1}

\begin{center}
COURBES ELLIPTIQUES :\quad BIBLIOGRAPHIE
\end{center}

1)\quad \uline{Fonctions elliptiques et modulaires}

C.JORDAN - Cours d'Analyse, Gauthi\makebox[0pt][l]{\deleted{\unsure{o|e}}}er-Villars

A.HURWITZ et R.COURANT - Vorlesungen über allgemeine Funktionentheorie und elliptische Funktionen, Springer (4ème édition - 1964)

R.FRICKE - Die elliptischen Funktionen und ihre Anwendungen, Teubner, Leipzig, 1922

2)\quad \uline{Multiplication complexe} (\uline{méthodes analytiques})

H.WEBER - Algebra, vol.III, Strasbourg, 1908

R.FUETER - Vorlesungen über die singulären Moduln und die komplexe Multiplik\unsure{.|a}tion der elliptischen Funktionen, I, 1924, II, 1927, Teubner.

H.HASSE - Neue Begründung der komplexen Multiplikation, I, Journ.Crelle, 157, p.115-139 (1927); II, \uline{ibid}., 165, p.64-88 (1931)

Seminar on complex multiplication, Inst.Adv.Study, Princeton, 1957-1958

M.DEURING - Die Klassenkörper der komplexen Multiplikation, Enz.Math.Wiss., I 2, 23, Teubner, 1958

3)\quad \uline{Corps de fonctions algébriques d'une variable}

C.CHEVALLEY - Introduction to the theory of algebraic functions of one variable, Math.Surveys, VI, New-York, 1951

M.EICHLER - Einführung in die Theorie der algebraischen Zahlen und Funktionen, Birkhäuser Verlag, 1963
% archive sheet number: 82
\clearpage
\scanpage{84}
\reversemarginpar
\margin{\figure{\tikz[x=1em,y=1em]{\draw[black!17,line width=.2pt] (0,0) -- (1.85,-.02); \fill (0,0) ellipse (.17 and .085); \fill (1.8,-.03) -- (1.86,.055) -- (1.94,.02) -- (1.9,-.06) -- cycle;}}}
\normalmarginpar
\mspage{B-2}

4)\quad \uline{Courbes algébriques}

\reversemarginpar
\margin{\handwritten{\figure{\tikz[x=1em,y=1em]{\draw[black!58,line width=.36pt] (0,-1.7) .. controls (.43,-1.1) and (1.08,-.39) .. (1.62,0);}}}}
A.WEIL - Sur \makebox[0pt][l]{\deleted{\unsure{-|=}}}les courbes algébriques et les variétés qui s'en déduisent, Hermann, 1948

J-P.SERRE - Groupes algébriques et corps de classes, Hermann, 1959

5)\quad \uline{Courbes elliptiques} (\uline{théorie algébrique})

H.HASSE - Zur Theorie der abstrakten elliptischen Funktionenkörper, I, II, III, Journ.Crelle, 175 (1936)

M.DEURING - Die Typen der Multiplikatorenringe elliptischer Funktionenkörper, Abh.Hamb., 16, 1949, p.32-47

6)\quad \uline{Schémas modulaires} (\uline{théorie algebrique})

J.IGUSA - Fibre systems of Jacobian varieties, III, Amer.J. of Maths., 81, 1959, p.453-476

J.IGUSA - On the transformation theory of elliptic functions, Amer.J.of Maths., 81, 1959, p.436-452

J.IGUSA - Kroneckerian model of fields of elliptic modular functions, Amer.J.of Maths., 81, 1959, p.561-577

7)\quad \uline{Variétés abéliennes} - \uline{Schémas abéliens}

\margin{\handwritten{\figure{\tikz[x=1em,y=1em]{\draw[black!54,line width=.33pt] (.42,-2.53) .. controls (.15,-2.64) and (.01,-2.62) .. (0,-2.58) .. controls (.68,-2.03) and (1.74,-1.4) .. (2.7,-.03);}}}}
A.WEIL - Variétés ab\unsure{e|é}liennes et courbes algébriques, Hermann, 1948

\margin{\handwritten{\figure{\tikz[x=1em,y=1em]{\draw[black!42,line width=.31pt] (0,-2.1) .. controls (.54,-1.55) and (1.12,-.64) .. (1.7,0);}}}}
S.LANG - Abelian Varieties, Interscience Tracts, n\textsuperscript{o}7, 1959

A.GROTHENDIECK - Fondements de la Géométrie algébrique (extraits du séminaire Bourbaki 1957-62), exp.232 et 236

A.NERON - Modèles minimaux des variétés abéliennes sur les corps locaux et globaux, Publ.Math.IHES, n\textsuperscript{o}21, 1964

D.MUMFORD - Geometric invariant theory, Ergebn.Math., 1965
\normalmarginpar
% archive sheet number: 83
\clearpage
\scanpage{85}
\reversemarginpar
\margin{\figure{\tikz[x=1em,y=1em]{\draw[black!18,line width=.24pt] (0,0) -- (1.81,-.05); \fill (0,.03) -- (.22,.08) -- (.23,-.08) -- (0,-.06) -- cycle; \fill[black!50] (.74,0) circle (.033); \fill[black!50] (1.16,-.035) circle (.033); \fill (1.68,-.14) -- (1.77,.035) -- (1.86,.05) -- (1.94,0) -- (1.82,-.05) -- cycle;}}}
\normalmarginpar
\mspage{B-3}

\makebox[0pt][l]{\deleted{\unsure{3|8}}}8)\quad \uline{Multiplication complexe} (\uline{théorie algébrique})

M.DEURING - Algebraische Begründung der komplexen Multiplikation, Abh.Hamb., 16, 1949, p.32-47

\reversemarginpar
\margin{\handwritten{\figure{\tikz[x=1em,y=1em]{\draw[black!47,line width=.34pt] (0,-2.04) .. controls (.73,-1.35) and (1.24,-.52) .. (1.83,0);}}}}
G.SHIMURA et Y.TAN\unsure{1|I}YAMA - Complex multiplication of abelian varieties, Publ.Math.Soc.Japan, n\textsuperscript{o}6, 1961

A.WEIL - On the theory of complex multiplication, Proc.Tokyo-Nikko Symp., p.9-22, 1955

M.DEURING - Die Zetafunktion einer algebraischen Kurve vom Geschlechte Eins, I, II, III, IV, Gött.Nach., 1953, 1955, 1956, 1957

9)\quad \uline{Autres questions arithmétiques relatives aux courbes elliptiques}

\margin{\handwritten{\figure{\tikz[x=1em,y=1em]{\draw[black!52,line width=.34pt] (0,-2.27) .. controls (.65,-1.92) and (1.98,-.64) .. (2.87,0);}}}}
S.LANG - Diophantine Geometry, Interscience Tracts, n\textsuperscript{o}11, 1962

J.CASSELS - Arithmetic on an elliptic curve, Proc.Int.Cong., 1962, p.234-246

J-P.SERRE - Groupes de Lie $\mathord{\rlap{\handwritten{$/$}}\unsure{\mathrm{l}|1}}$-adiques attachés aux courbes elliptiques, Colloque de Clermont, 1964

J.TATE - Rational points on elliptic cur\makebox[0pt][l]{\deleted{\unsure{:|v}}}\kern.11em\raisebox{.18ex}{v}es , Haverford College, 1961

10)\quad \uline{Multiplication complexe ``formelle''}

J.LUBIN - One parameter formal Lie groups over $\mathrm{p}$-adic integer rings, Ann.of Maths., 1965

J.LUBIN et J.TATE - Formal complex multiplication in local fields, Ann.of Maths., 1965

\begin{center}
***
\end{center}
\normalmarginpar
% archive sheet number: 84
